\section*{Chapitre \ref{ch:mx:the-almgren-chriss-framework} --- Le cadre d'Almgren--Chriss}

\begin{solution}{exo:mx:the-almgren-chriss-framework:1}
$\kappa=\sqrt{10^{-6}\times1/2.53\times10^{-7}}=\sqrt{3.95}=1.99$ par jour.
\end{solution}

\begin{solution}{exo:mx:the-almgren-chriss-framework:2}
$\E=\eta X^2/T=2.53\times10^{-7}\times10^{12}=253\,000$ dollars~; $\operatorname{Var}=\sigma^2X^2T/3$, soit un écart type de $10^6/\sqrt3=577\,000$ dollars
(576\,000 sur des intervalles d'une minute).
\end{solution}

\begin{solution}{exo:mx:the-almgren-chriss-framework:3}
5.9 des 6.5 heures. Sur une longue fenêtre, $x(t)\approx Xe^{-\kappa t}$, donc 95~\% sont vendus au bout de $\ln20/\kappa$.
\end{solution}

\begin{solution}{exo:mx:the-almgren-chriss-framework:4}
Le coût temporaire croît comme $\kappa$ (exécuter plus vite coûte davantage par action, sur chaque action), tandis que le risque évité en se pressant est faible
une fois l'ordre déjà vendu tôt~; être patient ajoute du risque lentement, car la variance est bornée par celle de la droite.
\end{solution}

\begin{solution}{exo:mx:the-almgren-chriss-framework:5}
Le glissement de chaque tranche ne mesure que le carnet qu'elle a consommé~; l'impact des tranches antérieures, qui n'a pas décru, fait monter les prix que paient les tranches suivantes, un
impact permanent (ou à décroissance lente) que le modèle, avec $\gamma=0$, laisse de côté.
\end{solution}

\begin{solution}{exo:mx:the-almgren-chriss-framework:6}
$\sigma$ a été estimé à partir de variations d'une minute et mis à l'échelle comme si les prix suivaient une marche aléatoire~; or le marché revient à la moyenne (les fondamentalistes ramènent
le prix), si bien que la variance sur trente minutes est inférieure à trente fois la variance d'une minute.
\end{solution}

\begin{solution}{exo:mx:the-almgren-chriss-framework:7}
L'écart moyen vaut $\gamma Q/2$~: $2\times(5.2-1.2)/20\,000=4.0\times10^{-4}$ tick par action pour $\kappa T=0$ et $2\times(7.4-3.6)/20\,000=3.8\times10^{-4}$ pour
$\kappa T=6$~: cohérent, comme doit l'être un terme permanent indépendant du calendrier.
\end{solution}

\begin{solution}{exo:mx:the-almgren-chriss-framework:8}
Ses transactions sont certaines~; son coût ne l'est pas~: il détient la position le plus longtemps, et le prix bouge pendant ce temps. Son risque est le plus élevé de la frontière,
soit un écart type de 576\,000~USD pour le bloc.
\end{solution}

\begin{solution}{pb:mx:the-almgren-chriss-framework:1}
\textbf{1.} Voir les définitions.
\textbf{2.} $\E=\tfrac12\gamma X^2+\eta\int v^2$, $\operatorname{Var}=\sigma^2\int x^2$.
\textbf{3.} Chaque action vendue déplace le prix de $\gamma$ au moment où elle est vendue~; le total vaut $\tfrac12\gamma X^2$.
\textbf{4.} Euler--Lagrange~: $\eta\ddot x=\lambda\sigma^2x$ avec $x(0)=X$, $x(T)=0$.
\textbf{5.} À une participation de 10~\%, la loi en racine carrée donne $0.8\sigma\sqrt{0.1}$ du prix~; en divisant par le taux, on obtient $\eta=2.53\times10^{-7}$ dollar
par action et par jour.
\textbf{6.} $\kappa=1.99$ par jour~; 300\,000, 470\,000 et 521\,000~USD.
\textbf{7.} Après 5.9 heures.
\textbf{8.} 253\,000, 576\,000 et 585\,000~USD.
\textbf{9.} Les couples (risque, coût) des calendriers optimaux, pour toutes les aversions au risque.
\textbf{10.} Patient~: 258\,000, 542\,000, 551\,000 (+5.8~\%)~; pressé~: 506\,000, 352\,000, 630\,000 (+20.9~\%).
\textbf{11.} Patience~: la frontière est plate autour de l'optimum vers une \omterm{def:mx:the-almgren-chriss-framework:urgency}{urgence} plus faible, et abrupte vers une \omterm{def:mx:the-almgren-chriss-framework:urgency}{urgence} plus forte.
\textbf{12.} Le nombre de dollars de coût espéré que le gestionnaire accepterait de payer pour retirer un dollar carré de variance~: un énoncé sur l'importance du risque de prix de la journée
compte pour le fonds.
\textbf{13.} 20\,000 actions sur 30 minutes, en 30 tranches, à des \omterm{def:mx:the-almgren-chriss-framework:urgency}{urgences} de 0, 2 et 6, huit sessions chacune.
\textbf{14.} Moyennes de 5.2, 5.3 et 7.4 ticks~; écarts types de 8.0, 6.5 et 3.8.
\textbf{15.} $\eta=0.11$, $\sigma=0.69$~; moyennes de 1.2, 1.4, 3.6 et écarts types de 16.4, 13.2, 7.6.
\textbf{16.} L'accumulation de l'impact d'une tranche à l'autre (une part permanente) et le retour à la moyenne du prix.
\textbf{17.} Environ $4\times10^{-4}$ tick par action.
\textbf{18.} \emph{Résultat nommé}~: avec $\lambda=10^{-6}$, le calendrier optimal vend 95~\% en 5.9 heures, espère 300\,000~USD avec un écart type de 470\,000~;
un calendrier deux fois trop patient coûte 5.8~\% d'équivalent certain de plus, et un calendrier deux fois trop pressé 20.9~\% de plus.
\textbf{19.} Le classement des calendriers est fiable~; les niveaux dépendent des calibrations, que la simulation montre pouvoir être décalées d'un facteur trois.
\textbf{20.} Le risque qui accompagne chaque coût, et donc le choix du coût qu'il vaut la peine de payer pour le réduire.
\end{solution}

\begin{solution}{iq:mx:the-almgren-chriss-framework:1}
Exécuter vite coûte l'impact temporaire~; exécuter lentement fait supporter le risque de prix. L'optimum concentre l'ordre en début de période~: $x(t)=X\sinh(\kappa(T-t))/\sinh(\kappa T)$, soit une
droite lorsque le risque n'importe pas et une décroissance exponentielle lorsqu'il importe beaucoup.
\iqlookfor{Le compromis~; la forme en sinh~; les cas limites.}
\end{solution}

\begin{solution}{iq:mx:the-almgren-chriss-framework:2}
Minimiser $\int(\eta\dot x^2+\lambda\sigma^2x^2)dt$~; équation d'Euler--Lagrange $\eta\ddot x=\lambda\sigma^2x$~; les conditions aux limites donnent le rapport en sinh.
\iqlookfor{Calcul des variations~; conditions aux limites.}
\end{solution}

\begin{solution}{iq:mx:the-almgren-chriss-framework:3}
$\kappa$ est multiplié par $\sqrt2$~: le calendrier concentre davantage en début de période, le coût espéré augmente et le risque diminue~; le client se déplace le long de la frontière vers l'exécution plus rapide.
\iqlookfor{Sens de chaque variation~; la racine carrée.}
\end{solution}

\begin{solution}{iq:mx:the-almgren-chriss-framework:4}
Temporaire~: régresser le glissement de chaque tranche par rapport au milieu, sur son taux d'exécution~; permanent~: la dérive du prix sur toute l'exécution, au-delà de la
part temporaire, ou le prix après la fin de l'exécution~; les deux sur de nombreux ordres, la loi en racine carrée servant de contrôle.
\iqlookfor{Estimations au niveau de la tranche et de l'ordre~; taille de l'échantillon.}
\end{solution}

\begin{solution}{iq:mx:the-almgren-chriss-framework:5}
Étant donné l'ordre et l'état du marché, renvoyer les positions cibles aux instants demandés (et un moyen de les mettre à jour à mesure que le marché évolue)~; la couche des \omterm{def:mx:the-almgren-chriss-framework:parent}{ordres
enfants} négocie vers ces cibles. Le rendre sans état, ou explicite sur son état, pour que les calendriers puissent être testés hors ligne.
\iqlookfor{Cibles dans le temps~; séparation avec le placement.}
\end{solution}

\begin{solution}{iq:mx:the-almgren-chriss-framework:6}
Impact non linéaire et transitoire, profils intrajournaliers de volume et de volatilité, signaux, fourchette et frais, adaptation au marché. Pour une action liquide, le
profil de volume et l'impact transitoire comptent en général le plus.
\iqlookfor{Une liste~; une hiérarchisation avec une justification.}
\end{solution}
